\chapter{JavaScript e il DOM}

\section{Il Document Object Model}

\begin{definizione}[DOM]
Il \textbf{DOM} (Document Object Model) \`e la struttura dati che rappresenta il documento
HTML e mantiene in memoria tutte le informazioni sulla sua visualizzazione. \`E un
\textbf{albero}: nodo \emph{root}, nodi intermedi (tipicamente elementi HTML) e foglie
(testo, commenti). Viene inizializzato da HTML e CSS, pu\`o essere modificato durante la vita
della pagina; la sua vita termina alla chiusura o al ricaricamento. Nell'ambiente-browser
\`e disponibile come oggetto \textbf{document}.
\end{definizione}

\section{Interfacce e la DOM API}

La parola \emph{interfaccia} denota il ``fronte esterno'' di un servizio, che permette ad
altri di utilizzarlo:
\begin{itemize}
  \item la \textbf{UI} (user interface) \`e il fronte verso gli \emph{utenti};
  \item l'\textbf{API} (application programming interface) \`e il fronte verso \emph{altri
        software / i programmatori}: dice come usare i servizi nel codice;
  \item l'\textbf{object interface} di un oggetto \`e costituita dalle sue propriet\`a e dalle
        \emph{signature} dei metodi: come l'oggetto pu\`o essere usato, a prescindere da come
        implementa il comportamento. \`E un'astrazione del concetto di classe (la classe
        contiene anche la componente \emph{operazionale}: come si costruiscono gli oggetti e
        come i metodi implementano il comportamento); una stessa object interface pu\`o essere
        implementata da classi diverse.
\end{itemize}
La specifica di un'interfaccia pu\`o essere \emph{a posteriori} (realizzo il servizio, poi ne
pubblico l'API) o \emph{a priori}, come \textbf{contratto}: l'elenco di funzionalit\`a \`e un
requisito che pu\`o essere implementato da fornitori diversi (analogamente, un mockup \`e una
specifica a priori di una UI).

\begin{esamebox}
\emph{``Illustrare le macrocategorie in cui \`e possibile suddividere le funzionalit\`a
offerte dalla DOM API del browser.''}
\end{esamebox}

\begin{definizione}[DOM API]
La \textbf{DOM API} \`e un'interfaccia del secondo tipo: una \textbf{specifica, stabilita dal
W3C}, di come deve essere il servizio di interrogazione e aggiornamento del DOM. Poich\'e il
servizio \`e offerto ai programmatori JavaScript \emph{dai browser}, sono i browser a giocare
il ruolo di fornitori e a doversi attenere alla specifica. \`E fatta di \emph{oggetti} che
devono essere forniti (il DOM stesso e i suoi elementi) e di \emph{object interfaces} che tali
oggetti devono rispettare; essendo ``solo'' una specifica \textbf{non contiene componenti
imperative, quindi non definisce classi} (sta a chi implementa decidere se usarle).
\end{definizione}

JavaScript non prevede le object interface come costrutti del linguaggio (servirebbero solo
al controllo statico dei tipi, che JS non ha): dire che i nodi DIV ``implementano
l'interfaccia \texttt{HTMLDivElement}'' significa solo che quegli oggetti hanno tutte le
propriet\`a e i metodi elencati nella specifica W3C. La gerarchia delle interfacce comprende,
fra le altre: \texttt{EventTarget} (radice), \texttt{Node}, \texttt{Window},
\texttt{Document}, \texttt{DocumentFragment}, \texttt{Element}, \texttt{HTMLElement} e le sue
specializzazioni (\texttt{HTMLDivElement}, \texttt{HTMLInputElement},
\texttt{HTMLParagraphElement}, \texttt{HTMLFormElement}, \texttt{HTMLImageElement}\dots),
\texttt{Attr}, \texttt{CharacterData} con \texttt{Text} e \texttt{Comment}.

\begin{importante}[Le tre macrocategorie della DOM API]
\begin{enumerate}
  \item \textbf{Navigazione / attraversamento del DOM}: trovare uno o pi\`u elementi in base
        alla struttura dell'albero e/o a tipo, classe, id, e leggerne le caratteristiche.
        Gestita per lo pi\`u da propriet\`a e metodi dell'interfaccia \texttt{Element}.
  \item \textbf{Modifica del DOM}: aggiungere, rimuovere e spostare elementi; modificare
        caratteristiche (classi, attributi), contenuti e stili. Anch'essa per lo pi\`u su
        \texttt{Element}.
  \item \textbf{Gestione degli eventi}: registrare (subscribe/unsubscribe) funzioni come
        \emph{listener} degli eventi UI di interesse, specifiche per tipologia di evento e per
        elemento sorgente; al verificarsi dell'evento il listener riceve le informazioni
        dettagliate. Gestita tramite l'interfaccia \texttt{EventTarget}.
\end{enumerate}
Lo schema tipico della programmazione UI in ``vanilla JavaScript'' \`e un loop su:
(1) trovare nel DOM gli elementi su cui intervenire, (2) registrarvi listener,
(3) al verificarsi dell'evento, modificare opportunamente il DOM.
\end{importante}

\section{Navigare il DOM}

\subsection{Ricerca per tipo, classe, id}
Definiti sia da \texttt{Element} (ricerca nel sottoalbero) sia da \texttt{Document} (ricerca
sull'intero documento):
\begin{itemize}
  \item \texttt{getElementById(id)} $\rightarrow$ \texttt{Element}: l'elemento con quell'id,
        o \texttt{null};
  \item \texttt{getElementsByTagName(tag)} $\rightarrow$ \texttt{HTMLCollection}
        (\emph{live});
  \item \texttt{getElementsByClassName(cls)} $\rightarrow$ \texttt{HTMLCollection}
        (\emph{live}).
\end{itemize}

\subsection{Attraversamento dell'albero}
\begin{itemize}
  \item da \texttt{Document}: \texttt{document.documentElement} (la radice),
        \texttt{document.head}, \texttt{document.body};
  \item da \texttt{Element}: \texttt{element.children} $\rightarrow$
        \texttt{HTMLCollection} (\emph{live}) dei figli che sono a loro volta elementi;
  \item da \texttt{Node}: \texttt{node.childNodes} $\rightarrow$ \texttt{NodeList}
        (\emph{live}) di \emph{tutti} i figli (elementi o meno);
        \texttt{node.parentElement} (il genitore \emph{se \`e un elemento}, altrimenti
        \texttt{null}); \texttt{node.parentNode} (il genitore di qualunque tipo, o
        \texttt{null}).
\end{itemize}

\subsection{Ricerca per selettore}
Usa i selettori CSS gi\`a visti. Discendente (da \texttt{Document} o da un \texttt{Element}):
\begin{itemize}
  \item \texttt{querySelector(sel)} $\rightarrow$ il \emph{primo} elemento corrispondente, o
        \texttt{null};
  \item \texttt{querySelectorAll(sel)} $\rightarrow$ \texttt{NodeList} \textbf{statica} di
        tutti i corrispondenti.
\end{itemize}
Ascendente (solo da un \texttt{Element}): \texttt{closest(sel)} $\rightarrow$ il primo
elemento corrispondente \emph{risalendo verso la radice}, o \texttt{null}.

\begin{attenzione}[HTMLCollection e NodeList: live vs statiche]
Sono collezioni \emph{array-like} ma \textbf{non array}: non supportano i metodi di
\texttt{Array}, ma permettono for classici e \texttt{for..of} (introdotte agli albori della
DOM API, mantenute per retrocompatibilit\`a). Le ricerche per tag/classe e le propriet\`a di
attraversamento restituiscono collezioni \textbf{live}: se il DOM cambia, i contenuti si
aggiornano. La ricerca per selettore restituisce collezioni \textbf{statiche}: per risultati
aggiornati bisogna \emph{ripetere la query}. Conversione:
\texttt{Array.from(coll)} -- ma l'array risultante \`e sempre e comunque statico.
\end{attenzione}

\section{Modificare il DOM}

\subsection{Contenuti (modifica bulk)}
\begin{itemize}
  \item \texttt{element.innerText} / \texttt{element.innerHTML} $\rightarrow$ string: il
        contenuto come testo o come HTML; \emph{modificandoli si sostituisce completamente il
        sottoalbero} che fa capo all'elemento.
\end{itemize}

\subsection{Creazione, rimozione, inserimento}
\begin{itemize}
  \item \texttt{document.createElement(tag)}: crea un \texttt{Element} \emph{non ancora nel
        DOM} (radice di un sottoalbero a parte);
  \item \texttt{element.remove()}: elimina l'elemento dal DOM;
  \item i tre metodi ``ortodossi'' di inserimento, in posizione contigua all'elemento su cui
        sono chiamati:
        \begin{itemize}
          \item \texttt{insertAdjacentElement(where, anElem)}: \texttt{anElem} pu\`o essere
                nuovo (da \texttt{createElement}) oppure gi\`a nel DOM -- in tal caso viene
                \textbf{spostato} (con tutto il suo sottoalbero);
          \item \texttt{insertAdjacentHTML(where, htmlString)}: costruisce un nuovo
                sottoalbero dalla stringa HTML e ne inserisce la radice;
          \item \texttt{insertAdjacentText(where, textString)}: crea e inserisce un nodo
                \texttt{Text}.
        \end{itemize}
        La posizione \texttt{where} \`e una stringa fra:
        \textbf{beforebegin} (fratello \emph{predecessore}: subito prima del tag di apertura),
        \textbf{afterbegin} (\emph{primo figlio}: subito dopo il tag di apertura),
        \textbf{beforeend} (\emph{ultimo figlio}: subito prima del tag di chiusura),
        \textbf{afterend} (fratello \emph{successore}: subito dopo il tag di chiusura).
\end{itemize}

\subsection{Classi e attributi}
\begin{itemize}
  \item \texttt{element.classList.add(c)} / \texttt{.remove(c)} / \texttt{.toggle(c)}:
        aggiungono/tolgono/alternano una classe fra quelle dell'attributo \texttt{class};
        \texttt{.has(c)} dice se l'elemento la possiede.
        \emph{Poich\'e cambiando le classi cambiano le regole di stile applicate, \`e un modo
        molto efficace di ottenere interattivit\`a}: far apparire/scomparire elementi,
        evidenziarli, spostarli.
  \item \texttt{element.getAttribute(nome)}: valore dell'attributo (come stringa);
  \item \texttt{element.setAttribute(nome, valore)}: assegna il valore all'attributo.
\end{itemize}
Nota: modificare il DOM (smontare/rimontare un elemento) \`e cosa diversa dal semplice
\texttt{display:none}.

\section{L'event loop asincrono single-thread}

L'ambiente di esecuzione JavaScript interagisce col mondo esterno tramite \textbf{eventi},
che possono essere:
\begin{itemize}
  \item \textbf{esogeni}: generati spontaneamente da un agente esterno (tipicamente l'utente);
        completamente imprevedibili -- non \`e dato dire se e quando avverranno;
  \item \textbf{endogeni}: conseguenza di una \emph{richiesta al mondo esterno} generata
        dall'applicazione stessa (es. una fetch); \`e ragionevolmente certo che si
        verificheranno, salvo errori.
\end{itemize}
Gli eventi \textbf{si mettono in coda}; perch\'e vengano gestiti, l'ambiente deve conoscere la
funzione che risponde all'evento:
\begin{itemize}
  \item \textbf{listener} per gli eventi esogeni: sta ``in ascolto'' del verificarsi
        dell'evento (non prevedibile) e risponde ogni volta che accade;
  \item \textbf{callback} per gli endogeni: a ogni avvio della richiesta si specifica la
        funzione da richiamare (call back) al completamento.
\end{itemize}
La logica \`e la stessa: istruzioni che \emph{associano a un tipo di evento una funzione} da
eseguire quando l'evento si verifica.

\begin{importante}[Single-thread e non bloccante]
JavaScript ``consuma'' la coda degli eventi con un \textbf{singolo thread}: gli eventi sono
gestiti in sequenza, nell'ordine di arrivo. Se la gestione di un evento intraprende
un'elaborazione lunga, \emph{blocca l'intera applicazione}: per questo le elaborazioni
time-consuming (richieste HTTP, scrittura nel DB del browser) sono gestite \textbf{in modo
asincrono} -- le istruzioni che le avviano non si bloccano in attesa del risultato, ma
richiedono una callback che verr\`a invocata a calcolo terminato.
(\`E in realt\`a possibile eseguire compiti specifici, che non toccano il DOM, in thread
paralleli detti \textbf{Web Workers}, ma il core dell'esecuzione resta come descritto.)
\end{importante}

\section{Gli eventi del DOM}

Gli oggetti del DOM ricevono eventi in corrispondenza delle interazioni supportate:
\begin{itemize}
  \item tutti gli elementi: dispositivi di puntamento -- \texttt{click},
        \texttt{mousedown}, \texttt{mouseup}, \texttt{mouseenter}, \texttt{mouseleave},
        \texttt{mouseover}, \texttt{mouseout}, \texttt{mousemove} (sui touch screen gli
        eventi touch si possono trattare a parte o assimilare al mouse);
  \item molti elementi di input: tastiera -- \texttt{focus}, \texttt{blur},
        \texttt{keydown}, \texttt{keyup};
  \item gli elementi form: interazioni di livello pi\`u alto -- \texttt{submit},
        \texttt{reset};
  \item \texttt{document}: \texttt{scroll}; \texttt{window}: \texttt{resize}.
\end{itemize}
L'elenco esaustivo \`e lunghissimo, ma per la maggior parte delle applicazioni bastano gli
eventi del mouse.

\subsection{Registrare un listener}
Tutti gli oggetti del DOM che ricevono eventi implementano:
\begin{lstlisting}[language=JavaScript]
domObj.addEventListener(nomeEvento, listener);
domObj.removeEventListener(nomeEvento, listener);
\end{lstlisting}
\texttt{nomeEvento} \`e una stringa (\texttt{"click"}, \texttt{"mouseup"},
\texttt{"keydown"}\dots); il \texttt{listener} \`e una funzione con un parametro opzionale e
nessun valore di ritorno, chiamata quando l'evento si verifica su \texttt{domObj}; riceve un
\emph{oggetto evento} con le informazioni dettagliate. La DOM API definisce le interfacce di
questi oggetti: \texttt{Event} e le sue specializzazioni \texttt{MouseEvent},
\texttt{KeyboardEvent}, \texttt{WindowEvent}, \texttt{FocusEvent}, eventi form\dots
Gli eventi UI sono per natura \emph{asincroni} e vanno gestiti tenendone conto.

\section{Bubbling e capturing}

\begin{definizione}[Bubbling]
Gli eventi del DOM si generano su un elemento preciso e da l\`i \textbf{risalgono l'albero
fino alla radice}: l'evento viene emesso sia dall'elemento sorgente sia da tutti gli elementi
sul percorso verso la radice. Nel listener, dall'oggetto evento si distinguono:
\begin{itemize}
  \item \texttt{evento.target}: l'elemento \emph{sorgente};
  \item \texttt{evento.currentTarget}: l'elemento che sta emettendo l'evento \emph{in quel
        momento} (presumibilmente quello a cui abbiamo attaccato il listener).
\end{itemize}
\end{definizione}

\textbf{Perch\'e il bubbling?} Se volessimo sapere quando l'utente clicca ``sul footer'',
senza bubbling riceveremmo il click solo sui punti del footer \emph{senza} elementi interni
(lo sfondo) -- raramente ci\`o che serve: ``click sul footer'' di solito significa ``footer
incluso ci\`o che contiene''. Il bubbling mette in pratica questo concetto.
Per \textbf{bloccare} la propagazione oltre un certo elemento (esempio tipico: il
\emph{glasspane} delle finestre modali), nel listener si usa
\texttt{evento.stopPropagation();}.

\begin{definizione}[Capturing]
La specifica stabilisce che, prima del bubbling, l'evento \emph{giunga} al target
\textbf{discendendo l'albero}: questa fase \`e il \textbf{capturing}. Di default \`e nascosta
(durante la discesa i listener non vengono invocati); si pu\`o chiedere l'invocazione in fase
di capturing con:
\begin{lstlisting}[language=JavaScript]
domObj.addEventListener(nomeEvento, listener, {capture: true});
\end{lstlisting}
Anche qui \`e possibile \texttt{stopPropagation()}: bloccando in capturing, l'evento
\emph{non arriva mai} al target e il bubbling non inizia -- utile per gestioni totalmente
customizzate di certi eventi.
\end{definizione}

\subsection{Effetti di default}
Alcuni eventi hanno effetti predefiniti: il click su un elemento A naviga altrove, il tasto
destro apre il menu contestuale, TAB passa al controllo successivo\dots\ Si impediscono nel
listener con \texttt{evento.preventDefault();}. Per bloccarli in tutta un'area
dell'applicazione conviene farlo \emph{in fase di capturing}.
